\section{Structural limits}\label{sec:limits}

The algorithmic results use four hypotheses: a product domain, an upper
bound $L$ on coordinate curvature, global quadratic growth at a unique
optimizer, and a correction that accounts for rounding every coordinate.
This section shows that each hypothesis does real work. The results are:
\begin{itemize}
\item a unique optimizer without a quantitative growth bound does not help:
  box QP with bag size three and a unique optimizer is NP-hard, with
  $\log\kappa=O(I)$ on the hard instances
  (Proposition~\ref{lim:prop:unique}); hence $\kappa$ cannot be replaced
  by $\log\kappa$ unless $\mathrm P=\mathrm{NP}$;
\item in the value-oracle model the exponent $p/2$ of $\kappa$ is optimal
  (Proposition~\ref{lim:prop:oracle});
\item exact piecewise-quadratic messages can need exponentially many pieces
  at bag size three and $\kappa\le80$ (Proposition~\ref{lim:prop:messages});
\item growth toward a nonunique optimal set does not bound grid sizes
  (Proposition~\ref{lim:prop:setgrowth});
\item local affine constraints make the problem NP-hard at bag size three
  and $\kappa=1$ (Proposition~\ref{lim:prop:constraints});
\item bag-local corrections are not valid lower bounds
  (Proposition~\ref{lim:prop:star});
\item local moment consistency of any finite order does not give an error
  bound in which large positive curvature is harmless
  (Proposition~\ref{lim:prop:moments}).
\end{itemize}
Del Pia and Khajavirad prove that continuous box QP is strongly NP-hard at
treewidth two, even with bounded integer coefficients, and that quartic
box problems on paths are strongly NP-hard \cite{DelPiaKhajavirad2026}.
Their instances need not have a unique optimizer, so their hardness does
not by itself show that $\kappa$ must enter the running time.
Remark~\ref{lim:rem:dk} shows that instances of their construction with a
unique optimizer can have exponentially large $\kappa$, consistently with
Theorem~\ref{thm:approx}.

\subsection{Conditioning cannot be dropped}\label{sec:limits-conditioning}

The following construction refines the simple weak NP-hardness reduction
in \cite[Remark~2]{DelPiaKhajavirad2026}: a stronger concave penalty puts
every minimizer at a vertex, a lexicographic linear term makes the minimizer
unique, and the growth constant is computed explicitly.

\begin{proposition}[Unique optimizers without conditioning]\label{lim:prop:unique}
Let $a_1,\ldots,a_n$ be positive integers, $n\ge2$, and let $T$ be an
integer with $0\le T\le U:=\sum_ia_i$. Put $M=\sum_ia_i^2$ and
$\varepsilon=1/(n2^{n+1})$. On the box
\[
x\in[0,1]^n,\qquad s=(s_1,\ldots,s_{n-1})\in[-U,2U]^{n-1},
\qquad s_0:=0,\quad s_n:=T,
\]
minimize
\[
F(x,s)=\sum_{i=1}^n(s_i-s_{i-1}-a_ix_i)^2
+M\sum_{i=1}^nx_i(1-x_i)+\varepsilon\sum_{i=1}^n2^{i-1}x_i .
\]
Then:
\begin{enumerate}
\item The bags $\{s_{i-1},s_i,x_i\}$, with the constants $s_0,s_n$ omitted
  and arranged as a path in the order $i=1,\ldots,n$, form a tree
  decomposition with $p\le3$. All data have binary length polynomial in
  that of $(a,T)$.
\item $F$ has a unique minimizer $z^*=(x^*,s^*)$, and $x^*\in\{0,1\}^n$.
  If some $v\in\{0,1\}^n$ satisfies $a^{\T}v=T$, then
  $a^{\T}x^*=T$ and $\OPT<1/(2n)$; otherwise $\OPT\ge1/n$.
\item $L=4$ is a valid upper coordinate-curvature bound and
  \[
  g=\frac{\varepsilon}{n\,(1+2(n-1)M)}
  \]
  is a valid growth constant. Hence
  $\kappa\le2^{n+3}n^2(1+2nM)$ and $\log_2\kappa=O(I)$.
\end{enumerate}
\end{proposition}

\begin{proof}
Part 1 is immediate: $s_i$ lies in the consecutive bags $i$ and $i+1$, each
$x_i$ in one bag, and each square and unary term lies in a bag.

\emph{Elimination of the states.} For $x\in[0,1]^n$ let
$e(x)=a^{\T}x-T$, $\sigma_0(x)=0$, and
\[
\sigma_k(x)=\sum_{j\le k}a_jx_j-\frac kne(x)\qquad(k=1,\ldots,n),
\]
so that $\sigma_n(x)=T$. For a feasible $s$ put $\delta_k=s_k-\sigma_k(x)$,
with $\delta_0=\delta_n=0$. Then
$s_i-s_{i-1}-a_ix_i=-e(x)/n+\delta_i-\delta_{i-1}$, and
$\sum_i(\delta_i-\delta_{i-1})=0$. Expanding the squares gives the identity
\begin{equation}\label{lim:eq:split}
F(x,s)=\Phi(x)+\sum_{i=1}^n(\delta_i-\delta_{i-1})^2,\qquad
\Phi(x)=\frac{e(x)^2}n+M\sum_ix_i(1-x_i)+\varepsilon\sum_i2^{i-1}x_i .
\end{equation}
Because $0\le\sum_{j\le k}a_jx_j\le U$ and $|e(x)|\le U$, the vector
$(\sigma_1(x),\ldots,\sigma_{n-1}(x))$ lies in $[-U,2U]^{n-1}$.
Since $\delta_k=\sum_{i\le k}(\delta_i-\delta_{i-1})$, the Cauchy--Schwarz
inequality gives $\delta_k^2\le k\sum_i(\delta_i-\delta_{i-1})^2$. Summing
over $k\le n-1$ gives the path inequality
\begin{equation}\label{lim:eq:path}
\sum_{i=1}^n(\delta_i-\delta_{i-1})^2\ge\frac2{n(n-1)}\sum_{k=1}^{n-1}
\delta_k^2\qquad(\delta_0=\delta_n=0).
\end{equation}

\emph{The reduced function.} Its Hessian is
$\nabla^2\Phi=\tfrac2naa^{\T}-2MI\preceq0$, since
$\norm a^2=M$. Thus $\Phi$ is concave. At a vertex $v$,
$\Phi(v)=e(v)^2/n+\varepsilon b(v)$ with $b(v)=\sum_i2^{i-1}v_i$, a
bijection of $\{0,1\}^n$ onto $\{0,\ldots,2^n-1\}$. If $\Phi(v)=\Phi(v')$,
then $(e(v)^2-e(v')^2)/n=\varepsilon(b(v')-b(v))$; the left side lies in
$n^{-1}\Z$ and the right side has absolute value less than
$\varepsilon2^n=1/(2n)$, so both vanish and $v=v'$. Let $x^*$ be the unique
minimizing vertex, $\Phi^*=\Phi(x^*)$, and
$\gamma=\min_{v\ne x^*}(\Phi(v)-\Phi^*)$. We claim $\gamma\ge\varepsilon$.
If $e(v)^2=e(x^*)^2$, the difference is a positive integer multiple of
$\varepsilon$. Otherwise $e(v)^2-e(x^*)^2$ is a nonzero integer. It cannot be
negative, since then $\Phi(v)-\Phi^*\le-1/n+\varepsilon(2^n-1)<0$. Hence it
is at least one and $\Phi(v)-\Phi^*\ge1/n-\varepsilon(2^n-1)>1/(2n)\ge
\varepsilon$.

For $x\in[0,1]^n$ let $Y$ have independent coordinates with
$\Pr(Y_i=1)=x_i$. By concavity and Jensen's inequality,
$\Phi(x)=\Phi(\E Y)\ge\E\Phi(Y)\ge\Phi^*+\gamma\Pr(Y\ne x^*)$. Since $x^*$ is
binary, $\Pr(Y_i\ne x_i^*)=|x_i-x_i^*|\le1$, so
$\Pr(Y\ne x^*)\ge\max_i|x_i-x_i^*|\ge\frac1n\norm{x-x^*}^2$. Therefore
\begin{equation}\label{lim:eq:phigrowth}
\Phi(x)-\Phi^*\ge\frac\varepsilon n\norm{x-x^*}^2\qquad(x\in[0,1]^n).
\end{equation}

\emph{Uniqueness and growth.} By \eqref{lim:eq:split} and
\eqref{lim:eq:phigrowth}, $z^*=(x^*,\sigma(x^*))$ is the unique minimizer
and $\OPT=\Phi^*$. Let $z=(x,s)$ be feasible and $h=x-x^*$. Then
$\sigma_k(x)-\sigma_k(x^*)=\sum_jc_{kj}a_jh_j$ with
$c_{kj}=[j\le k]-k/n\in[-1,1]$, so
$\norm{\sigma(x)-\sigma(x^*)}^2\le(n-1)M\norm h^2$ and
\[
\norm{z-z^*}^2\le\norm h^2+2\norm\delta^2+2\norm{\sigma(x)-\sigma(x^*)}^2
\le(1+2(n-1)M)\norm h^2+2\norm\delta^2 .
\]
By \eqref{lim:eq:split}--\eqref{lim:eq:phigrowth},
$F(z)-\OPT\ge(\varepsilon/n)\norm h^2+\frac2{n(n-1)}\norm\delta^2$. Since
$g(1+2(n-1)M)=\varepsilon/n$ and $g\le\varepsilon/n<1/(n(n-1))$, this gives
$F(z)-\OPT\ge g\norm{z-z^*}^2$. On coordinate lines,
$\partial^2F/\partial s_k^2=4$ and $\partial^2F/\partial x_i^2=2a_i^2-2M\le0$,
so $L=4$ is valid and $\kappa\le4/g\le2^{n+3}n^2(1+2nM)$. Both $n$ and
$\log M$ are $O(I)$.

\emph{Decision.} If $a^{\T}v=T$ for a vertex $v$, then
$\OPT\le\Phi(v)\le\varepsilon(2^n-1)<1/(2n)$. As shown above,
$e(x^*)^2\le e(v)^2=0$. If no vertex solves the instance, every vertex has
$e(v)^2\ge1$, and $\OPT=\Phi(x^*)\ge1/n$.
\end{proof}

\begin{corollary}[The dependence on $\kappa$ is not polylogarithmic]
\label{lim:cor:nopolylog}
Unless $\mathrm P=\mathrm{NP}$, no algorithm computes, for every rational
continuous box QP with a supplied tree decomposition of maximum bag size at
most three and a unique minimizer, the optimal value within $2^{-q}$ in time
polynomial in $I+q+\log\kappa$. In particular, no algorithm achieves the
guarantee of Theorem~\ref{thm:approx} with $f(3,\kappa)$ bounded by a
polynomial in $\log\kappa$.
\end{corollary}

\begin{proof}
Apply such an algorithm to the instances of
Proposition~\ref{lim:prop:unique} with $q=\lceil\log_2(4n)\rceil$. An
approximation within $1/(4n)$ separates $\OPT<1/(2n)$ from $\OPT\ge1/n$,
and $\log\kappa=O(I)$. This would decide Subset Sum, which is NP-complete
\cite{GareyJohnson1979}, in polynomial time.
\end{proof}

Theorem~\ref{thm:approx} bounds $f(p,\kappa)$, for fixed $p$, by a
polynomial in $\kappa$: the state cap is $O(\sqrt\kappa\log(n+2))$ per
coordinate, and the arithmetic is polynomial in its numerator lengths.
Hence for fixed $p$ the unique-optimum problem is polynomial on instances
with $\kappa\le I^{O(1)}$ (Theorem~\ref{thm:exact}), while
Proposition~\ref{lim:prop:unique} makes it NP-hard on instances with
$\kappa\le2^{O(I)}$. The parameter $\kappa$ therefore locates the
complexity boundary at fixed width.


\begin{remark}[Conditioning of the bounded-coefficient reduction]
\label{lim:rem:dk}
The strongly NP-hard instances of \cite[Theorem~3]{DelPiaKhajavirad2026}
have bags of size three and coefficients bounded by absolute constants,
but their minimizers need not be unique. Consider their construction for
the two-item instance $a=(B,B-1)$, $T=B$, $B\ge3$, with objective $\Psi$
(their (20)--(24), all weights one) on the unit box. Here
$\ell=\lceil\log_2(2B)\rceil$, $D=2^\ell\ge2B$, and there are $5\ell+2$
variables. Every term of $\Psi$ is nonnegative. A zero of $\Psi$ encodes a
binary solution, and $(1,0)$ is the only one; all copied and state variables
are then determined by the recurrences. Hence $\Psi$ has a unique minimizer
$v^*$ with value zero. By the qualitative growth lemma for quadratics with a
unique minimizer, some $g>0$ is valid.

Encode instead the choice $(0,1)$ with every copy, digit and running-sum
recurrence exact. Only the final square
$(s_2-w_\ell)^2=((B-1-B)/D)^2$ is nonzero, so $\Psi(y')=D^{-2}$. All
$\ell$ copies of both choice bits change by one, so
$\norm{y'-v^*}^2\ge2\ell$ and every valid growth constant satisfies
$g\le1/(2\ell D^2)$. The largest diagonal entry of $\nabla^2\Psi$ is ten,
attained by the digit-state variables. Therefore
\[
\kappa\ge20\ell D^2\ge80\ell B^2 .
\]
For $B=2^m$ there are $5m+7$ variables and coefficients bounded by five, so
$\kappa$ is exponential in the input length, consistent with
Corollary~\ref{lim:cor:nopolylog}.

Negative curvature is also large relative to growth. Encode the fractional
choice $(1/B,1)$; since $B\cdot B^{-1}+(B-1)\cdot1=B$, every square
residual can be kept at zero with all variables in $[0,1]$, and
$\Psi(y)=q/B$ with $q=1-1/B$. The displacement $d=y-v^*$ annihilates every
affine residual, so $d^{\T}\nabla^2\Psi\,d=-2(q^2+1)$, coming from the
two endpoint penalties. With $R=\norm d^2$, every valid growth constant has
$g\le q/(BR)$, while $\nu\ge2(q^2+1)/R$. Hence
\[
\frac\nu g\ge2B\Bigl(q+\frac1q\Bigr)\ge4B .
\]
This bound on $\nu/g$ also holds for arbitrary positive weights on the
squares and a common positive weight on the two endpoint penalties: both
estimates scale with the penalty weight, and the square weights do not
enter. The reduction therefore does not contradict an algorithm
parameterized by $\nu/g$ either.
\end{remark}

\subsection{The exponent of $\kappa$ in the value-oracle model}
\label{sec:limits-oracle}

The corrected-grid lower bound needs only exact factor values. In a model
where factors are available only through evaluations, the exponent $p/2$ of
$\kappa$ cannot be improved. The construction is the classical hidden-well
argument \cite{NemirovskiYudin1983}, adapted to the present hypotheses.

\begin{proposition}[Value-oracle lower bound]\label{lim:prop:oracle}
Let $p\ge1$, $L>0$, $\kappa\ge8p$, $g=L/\kappa$, and $X=[0,1]^p$. For
$c\in X$ put
\[
F_c(x)=\min\Bigl\{gp,\ \frac L2\norm{x-c}^2\Bigr\}.
\]
Each $F_c$ has upper coordinate curvature $L$, the unique minimizer $c$
with $F_c(c)=0$, and growth $F_c(x)\ge g\norm{x-c}^2$ on $X$; a single
bag of size $p$ is a valid decomposition. Suppose a deterministic algorithm
receives $p$, $L$, $\kappa$ and value access to $F_c$ for an unknown $c$,
and always returns $y\in X$ with $F_c(y)<Lp/\kappa$. Then on some $F_c$ it
makes at least $(\kappa/(8p))^{p/2}-1$ evaluations. A randomized algorithm
that makes at most $Q$ evaluations succeeds, for some $c$, with probability
at most $(Q+1)(8p/\kappa)^{p/2}$.
\end{proposition}

\begin{proof}
On a coordinate line, $\frac L2\norm{x-c}^2-\frac L2x_i^2$ is affine in
$x_i$ and $gp-\frac L2x_i^2$ is concave; their minimum is concave. Thus
$t\mapsto F_c-Lt^2/2$ is concave on every coordinate line. For $x\in X$,
$\norm{x-c}^2\le p$, so $gp\ge g\norm{x-c}^2$; and $L/2\ge g$ because
$\kappa\ge2$. This proves growth, and $c$ is the unique zero.

Moreover $F_c(x)=gp$ unless $\norm{x-c}<r$, where $r^2=2gp/L=2p/\kappa$.
Since $\kappa\ge8p$, $2r\le1$. Let $\mathcal C$ consist of the points of
$X$ whose coordinates lie in $\{0,2r,4r,\ldots\}$. Then
$|\mathcal C|=(\lfloor1/(2r)\rfloor+1)^p\ge(2r)^{-p}=(\kappa/(8p))^{p/2}$,
and the open balls $B(c,r)$, $c\in\mathcal C$, are pairwise disjoint.

Run the deterministic algorithm with every answer equal to $gp$. Let
$x^1,\ldots,x^Q$ be its evaluation points and $y$ its output. Each of these
$Q+1$ points lies in at most one ball. If $Q+1<|\mathcal C|$, some
$c\in\mathcal C$ has a ball containing none of them. For this $F_c$ all
answers are $gp$, so the algorithm makes the same evaluations and returns
$y$, with $F_c(y)=gp=Lp/\kappa$. This contradicts the guarantee, so
$Q\ge|\mathcal C|-1$.

For a randomized algorithm, choose $c$ uniformly in $\mathcal C$ and fix
the random bits. The run on the constant answers has at most $Q$
evaluations and one output. If $F_c$ differs from $gp$ at none of the
evaluation points, the run on $F_c$ is the same and succeeds only if the
output lies in $B(c,r)$. Hence at most $Q+1$ choices of $c$ succeed.
Averaging over the random bits gives success probability at most
$(Q+1)/|\mathcal C|$ for a uniform $c$, hence for some $c$.
\end{proof}

\begin{remark}[Comparison with the algorithm, and the dependence on $p$]
\label{lim:rem:oracle}
With one bag ($n=p$, $N=1$), each stage of the capped algorithm evaluates
$F$ on at most $K^p$ grid points. By Lemma~\ref{lem:states} and the
trial schedule, $K=O(\sqrt\kappa\log(p+2))$ in the successful trial, and
the earlier trials cost no more. Thus the exponent $p/2$ of $\kappa$ is
attained, and Proposition~\ref{lim:prop:oracle} shows that it cannot be
improved in this model; the remaining gap is a factor
$(C\sqrt p\log(p+2))^p$ times the number of stages. The lower bound is also
exponential in $p$ once $\kappa\ge32p$.

In the explicit model, dependence on $p$ must be superpolynomial even at
$\kappa=1$, unless $\mathrm P=\mathrm{NP}$. For a graph $(V,E)$ with
$V=\{1,\ldots,n\}$ let
$F(x)=\sum_{ij\in E}(2x_ix_j-x_i-x_j)+2^{-n-1}\sum_i2^{i-1}x_i$ on
$[0,1]^n$. It is multilinear, so every $L>0$ is a valid curvature bound.
At vertices it equals minus the cut size plus a tie-breaking term below
one half; the minimizing vertex $x^*$ is unique and encodes a maximum cut.
For multilinear $F$ and the independent rounding $Y$ of $x$,
$F(x)=\E F(Y)$, so the argument leading to \eqref{lim:eq:phigrowth} gives
growth $g=2^{-n-1}/n$. Choosing $L=g$ gives $\kappa=1$ with one bag of size
$n$, and Max-Cut is NP-hard \cite{GareyJohnson1979}. Whether a white-box
lower bound of the form $\kappa^{\Omega(p)}$ holds is open.
\end{remark}

\subsection{Exact messages can be exponentially large}
\label{sec:limits-messages}

For forests, exact value functions of box QP are piecewise quadratic with
linearly many pieces \cite{DelPiaKhajavirad2026}. At bag size three, exact
messages can be exponentially large even when $\kappa$ is bounded and the
optimizer is unique. This is why the algorithm uses grids rather than exact
piecewise-quadratic messages.

\begin{proposition}[Exponential messages at bounded conditioning]
\label{lim:prop:messages}
For $m\ge2$ let $S_t\in[0,2^t-1]$ and $z_t\in[0,1]$, $t=1,\ldots,m$, put
$S_0:=0$, $r_t=S_t-2S_{t-1}-z_t$, and
\[
G_m(S,z)=S_m^2+\sum_{t=1}^mr_t^2+\frac18\sum_{t=1}^mz_t(1-z_t).
\]
Then:
\begin{enumerate}
\item the bags $\{S_{t-1},S_t,z_t\}$ ($S_0$ omitted) form a path
  decomposition with $p=3$, and $I=O(m^2)$;
\item $G_m(S,z)\ge\frac18(\norm S^2+\norm z^2)$ on the box, so zero is the
  unique minimizer and $g=1/8$ is valid;
\item $L=10$ is the exact maximum coordinate curvature, so $\kappa\le80$;
  moreover $\nabla^2G_m\succeq-\frac14I$;
\item let $\mu(v)=\min\{G_m(S,z):S_m=v\}$ for $v\in[0,2^m-1]$, the
  min-marginal of $S_m$. Up to the term $v^2$, it is the message sent
  across the separator $\{S_m\}$ to a leaf bag $\{S_m\}$ carrying $S_m^2$.
  If $[0,2^m-1]$ is covered by
  $K$ intervals on each of which $\mu$ coincides with a quadratic
  polynomial, or if $\mu=\min_{k\le K}q_k$ for quadratic polynomials $q_k$,
  then $K\ge2^{m-1}$.
\end{enumerate}
By Theorem~\ref{thm:approx}, the corrected-grid algorithm nevertheless
solves this family to accuracy $2^{-q}$ in time polynomial in $m+q$.
\end{proposition}

\begin{proof}
Part 1: $S_t$ lies in bags $t$ and $t+1$, and the endpoints $2^t-1$ have
$t$ bits.

Part 2. On the box $S_t\ge0$ and $z_t\ge0$, so
$S_{t-1}=(S_t-z_t-r_t)/2\le(S_t+|r_t|)/2$. By induction,
\[
0\le S_t\le2^{t-m}S_m+\sum_{j=t+1}^m2^{t-j}|r_j| .
\]
The vector $(2^{t-m})_{t\le m}$ has squared norm at most $4/3$. The map
$y\mapsto(\sum_{j>t}2^{t-j}y_j)_t$ equals $\sum_{l\ge1}2^{-l}J^l$, where $J$
shifts a vector by one place, so its operator norm is at most one. Hence $\norm S\le(2/\sqrt3)S_m+\norm r$ and
$\norm S^2\le\frac83S_m^2+2\norm r^2$. Also
$0\le z_t=S_t-2S_{t-1}-r_t\le S_t+|r_t|$, so
$\norm z^2\le2\norm S^2+2\norm r^2$. Thus
$\norm S^2+\norm z^2\le3\norm S^2+2\norm r^2\le8S_m^2+8\norm r^2\le8G_m$.

Part 3. The second derivatives are $2+8=10$ for $S_t$, $t<m$; $2+2=4$ for
$S_m$; and $2-\frac14$ for $z_t$. Writing $G_m$ as a sum of squares plus
$-\frac18\sum z_t^2$ plus linear terms gives $\nabla^2G_m+\frac14P_z\succeq0$,
where $P_z$ projects onto the controls.

Part 4. Let $W(v)=\mu(v)-v^2
=\min\{\sum_tr_t^2+\frac18\sum_tz_t(1-z_t):S_m=v\}\ge0$. The minimum is
attained by compactness. Hence $W(v)=0$ exactly when some feasible point has
$S_m=v$, all $r_t=0$ and binary $z$. Then $S_t=2S_{t-1}+z_t$ and
$v=\sum_t2^{m-t}z_t\in\{0,\ldots,2^m-1\}$. Conversely, for an integer $j$ in
this range, its binary digits give $S_t=\lfloor j/2^{m-t}\rfloor\in
[0,2^t-1]$. So $W$ vanishes exactly at the $2^m$ integers of
$[0,2^m-1]$.

In the first representation, on an interval with nonempty interior,
$W$ coincides with a polynomial of degree at most two. It is not
identically zero, since $W>0$ off the integers. So it has at most two zeros,
and the interval contains at most two integers. A degenerate interval
contains at most one. Hence $2K\ge2^m$. In the second representation, each
integer $j$ has some $k$ with $q_k(j)=\mu(j)$. Then $q_k-v^2\ge W\ge0$ on the
interval and vanishes at $j$. Since $q_k-v^2\not\equiv0$, it vanishes at at
most two integers, and again $2K\ge2^m$.
\end{proof}

\subsection{Point growth cannot be replaced by set growth}
\label{sec:limits-setgrowth}

\begin{proposition}[Set growth does not bound grids]\label{lim:prop:setgrowth}
For $n\ge2$ let $F(x)=\sum_{i=1}^{n-1}(x_i-x_{i+1})^2$ on $[0,1]^n$, with
the path decomposition of bags $\{x_i,x_{i+1}\}$. The optimal set is the
diagonal $S=\{t\mathbf1:t\in[0,1]\}$, and
$F(x)\ge\frac2{n(n-1)}\dist(x,S)^2$; for $n=2$, $F=2\dist(x,S)^2$. Every
valid $L$ is at least $2$. For any grids $G_i\ni0,1$ and the corrections
$d_i(v)=Lw_i(v)^2/8$,
\[
\beta=\min_{y\in\prod_iG_i}(F-D)(y)\le-\frac L8W_j^2
\qquad\hbox{for every }j,
\]
where $W_j$ is the largest gap of $G_j$. The filtering rule of
Proposition~\ref{prop:filter} never removes an interval. Hence every
corrected-grid certificate of accuracy $\varepsilon$, including its
filtering history, has a last-stage grid with at least
$1+1/(2\sqrt\varepsilon)$ nodes in every coordinate. Its size is not
polynomial in $\log(1/\varepsilon)$, although $L$ divided by the set-growth
constant is at most $2n(n-1)$.
\end{proposition}

\begin{proof}
For $x\in[0,1]^n$, $\dist(x,S)^2\le\sum_i(x_i-x_1)^2$, and
$(x_i-x_1)^2\le(i-1)F(x)$ by Cauchy--Schwarz, which gives the set-growth
bound. For $n=2$, $\dist(x,S)^2=(x_1-x_2)^2/2$. The coordinate second
derivatives are $2$ or $4$, so $L\ge2$.

Fix $j$, and let $v_j\in G_j$ be an endpoint of a largest gap, so
$w_j(v_j)=W_j$. Choose the remaining coordinates outward along the path. If
$v_k$ has been chosen and $i=k\pm1$ is the next index, let $[\alpha,\alpha']$
be a gap of $G_i$ containing $v_k$, of length $\gamma_i$, and let $v_i$ be
its endpoint nearest to $v_k$. Then $(v_k-v_i)^2\le\gamma_i^2/4$ and
$w_i(v_i)\ge\gamma_i$. Every edge of the path is used exactly once, so
\[
F(v)-D(v)\le\sum_{i\ne j}\frac{\gamma_i^2}4-\frac L8W_j^2
-\sum_{i\ne j}\frac L8\gamma_i^2\le-\frac L8W_j^2,
\]
because $L\ge2$.

Every $t\in[0,1]$ is a coordinate of the optimal point $t\mathbf1$. By
Proposition~\ref{prop:cellwise}, every interval $[a,b']$ of every grid satisfies
$\min\{m_i(a),m_i(b')\}\le F(t\mathbf1)=0\le U$ for $t\in[a,b']$, so the
filtering rule keeps it. Thus every stage covers $[0,1]^n$, and the final
gap is at least $U-\beta\ge LW_j^2/8\ge W_j^2/4$ since $U\ge\OPT=0$. A gap
at most $\varepsilon$ forces $W_j\le2\sqrt\varepsilon$, hence
$|G_j|\ge1+1/(2\sqrt\varepsilon)$.
\end{proof}

For finitely many optimal values per coordinate, the union-of-cells variant
of the TU section remains valid. The proposition concerns a continuum of
optimal coordinate values. It is a limit of corrected-grid certificates, not
of the problem: here $F$ is a sum of squares.

\subsection{Product domains are essential}\label{sec:limits-constraints}

\begin{proposition}[Local affine constraints]\label{lim:prop:constraints}
Let $a_1,\ldots,a_n$ and $B$ be positive integers with $a_i\le B$, and put
$a_0=B$. Consider binary $x_0,\ldots,x_n$, continuous
$y_0,\ldots,y_{n+1}\in[0,1]$, and the constraints
\[
y_0=0,\qquad y_{n+1}=1,\qquad y_{i+1}=y_i+\frac{a_i}Bx_i\quad(i=0,\ldots,n),
\]
with objective $F=2^{n+1}x_0+\sum_{i=1}^n2^{i-1}x_i$. Each constraint lies in
the bag $\{y_i,x_i,y_{i+1}\}$, and these bags form a path with $p=3$. The
feasible set is nonempty, the minimizer $z^*$ is unique, and
$F(z)-\OPT\ge\norm{z-z^*}^2/(2n+3)$ for every feasible $z$. Since $F$ is
affine, every $L>0$ is valid; with $L=1/(2n+3)$, $\kappa=1$. Moreover
$x_0^*=0$ if and only if some subset of $\{a_1,\ldots,a_n\}$ sums to $B$.
Consequently, unless $\mathrm P=\mathrm{NP}$, no algorithm solves such
problems in time $f(p,\kappa)\poly(I)$ for any function $f$, even when the
optimal value is required only within $1/2$.
\end{proposition}

\begin{proof}
Summing the recurrences gives $\sum_{i=0}^na_ix_i=B$. The choice $x_0=1$,
$x_i=0$ ($i\ge1$) is feasible. Since all $a_i>0$, it is the only feasible
choice with $x_0=1$. A choice with $x_0=0$ is feasible exactly when its
items sum to $B$: its partial sums then lie in $[0,1]$. For each binary $x$
the constraints determine $y$. Distinct feasible points therefore have
distinct binary vectors and distinct integer objective values. The dummy
value $2^{n+1}$ exceeds every value $\le2^n-1$ of a genuine subset. Hence
the minimizer is unique, and $x_0^*=0$ exactly for yes-instances. All
$2n+3$ coordinates lie in $[0,1]$, so $\norm{z-z^*}^2\le2n+3$, while
$F(z)-\OPT\ge1$ for $z\ne z^*$. An approximation within $1/2$ separates
$\OPT\le2^n-1$ from $\OPT=2^{n+1}$. The data have polynomial binary length,
and Subset Sum is NP-complete \cite{GareyJohnson1979}.
\end{proof}

In the product setting an affine objective is solved exactly by endpoint
dynamic programming. The difficulty here comes entirely from the
constraints. Independent coordinate rounding, the basis of
Proposition~\ref{prop:cellwise}, does not preserve the running-sum equations.

\subsection{The correction must be global}\label{sec:limits-star}

The correction $D$ charges for rounding every coordinate, including
coordinates outside a bag. A tempting alternative charges a bag cell only for
its own coordinates. The next example shows that this alternative is not
valid.

\begin{proposition}[Bag-local corrections fail]\label{lim:prop:star}
Let $0<\varepsilon\le1/16$ and let $m>16$ be an integer. On $[0,1]^{m+1}$
put
\[
F(x,y)=x^2-\frac x2\sum_{i=1}^my_i+\sum_{i=1}^my_i^2
+\Bigl(\frac m{16}-1-\varepsilon\Bigr)x,
\]
with the star decomposition of bags $\{x,y_i\}$. Then $L=2$ is valid, the
unique minimizer is $(1,\tfrac14,\ldots,\tfrac14)$, and $\OPT=-\varepsilon$.
Use the grid $\{0,\tfrac12,1\}$ in every coordinate, with uncorrected grid
min-marginals
$M(\xi,\eta)=\min\{F(x,y):\hbox{grid},\ x=\xi,\ y_1=\eta\}$ and grid
incumbent $U$. Consider any rule that discards the cell
$[\tfrac12,1]\times[0,\tfrac12]$ of the bag $\{x,y_1\}$ when every corner
satisfies $M(\xi,\eta)-e>U$. If $m>32e+8+16\varepsilon$, the rule discards
this cell, although it contains the bag projection $(1,\tfrac14)$ of the
optimizer. In particular no correction $e$ depending only on $|B|=2$,
$L=2$ and $h=1/2$ is valid; $e=|B|Lh^2/8=1/8$ fails for every $m\ge17$.
\end{proposition}

\begin{proof}
The coordinate second derivatives are $2$, so $L=2$ is valid. For fixed $x$,
each $y_i^2-xy_i/2$ is strictly convex with minimizer $x/4\in[0,1]$.
Eliminating the leaves gives $V(x)=(1-\frac m{16})x^2+(\frac m{16}-1-
\varepsilon)x$. This is strictly concave for $m>16$, with $V(0)=0$ and
$V(1)=-\varepsilon$, so the unique minimizer is $x=1$, $y_i=1/4$.

On the grid, $\min_{y\in\{0,1/2,1\}}(y^2-\xi y/2)=0$ for $\xi\in\{0,\tfrac12,
1\}$. Hence $M(\xi,\eta)=\xi^2+c\xi+\eta^2-\xi\eta/2$ with
$c=\frac m{16}-1-\varepsilon$. The four corners have values
$\frac m{32}-\frac14-\frac\varepsilon2$, this value plus $\frac18$, and twice
$\frac m{16}-\varepsilon$. The smallest is
$\frac m{32}-\frac14-\frac\varepsilon2$. The grid values at
$x\in\{0,\tfrac12,1\}$ are $0$,
$\frac m{32}-\frac14-\frac\varepsilon2>0$ and $\frac m{16}-\varepsilon>0$,
so $U=0$. The rule discards the cell exactly when
$\frac m{32}-\frac14-\frac\varepsilon2-e>0$. For $e=1/8$ and $m\ge17$ this
holds.
\end{proof}

The missing error is $\min_{\rm grid}F(x,\cdot)-V(x)=m\dist(x/4,\{0,\tfrac12,
1\})^2$. It comes from the outside coordinates, grows linearly in $m$, and
depends on the separator value, so message normalization does not remove it.
The global correction of Proposition~\ref{prop:cellwise} includes these coordinates.
Conditional recourse with certified outside bounds is the valid local
alternative.

